% !TeX spellcheck = fr_FR
\documentclass[french]{yLectureNote}

\title{Électromagnétisme 2 PS}
\subtitle{Physique}
\author{Paulhenry Saux}
\date{\today}
\yLanguage{Français}

\usepackage[utf8]{inputenc}
\usepackage[T1]{fontenc}
\usepackage{amsmath, amssymb, amsthm}
\usepackage{awesomebox}

\newcommand{\Lim}[1]{\lim\limits_{\substack{#1}}\:}
\renewcommand{\vec}{\overrightarrow}
\newcommand{\N}[0]{\mathbb{N}}
\newcommand{\dd}{\mathrm{d}}
\newcommand{\norm}[1]{||\vec{#1}||}
\newcommand{\fo}{\psi(\vec{r},t)}
\newcommand{\foe}{\psi(\vec{r},t)\*}
\newcommand{\HH}{\hat{H}}
\newcommand{\hb}{\hbar}
\newcommand{\lap}{\nabla^2}
\newcommand{\lapcc}{\frac{\partial^2 }{\partial x^2}+\frac{\partial^2 }{\partial y^2}+\frac{\partial^2 }{\partial z^2}}
\newcommand{\E}{\vec{E}(M)}
\newcommand{\B}{\vec{B}(M)}
\newcommand{\V}{V(M)}
\newcommand{\er}{\vec{e_{\rho}}}
\newcommand{\ep}{\vec{e_{\varphi}}}
\newcommand{\pip}{\pi^+}
\newcommand{\pim}{\pi^-}
\newcommand{\mv}{\mathcal{V}}
\DeclareMathOperator{\Div}{Div}
\newcommand{\rot}{\vec{\mathrm{rot}}}
\newcommand{\grad}{\vec{\mathrm{grad}}}

\begin{document}
\setcounter{chapter}{2}
\chapter{Équations de Maxwell}

\section{Conservation de la charge}

La charge électrique est une grandeur conservative. Toute variation de la charge dans un volume $V$ est due à un flux de courant à travers sa surface frontière $S$.

\begin{theorem}[Équation de conservation de la charge]
    En tout point $M$ à l'instant $t$, la densité volumique de charge $\rho$ et le vecteur densité de courant $\vec{j}$ sont liés par la forme locale :
    $$\Div \vec{j} + \frac{\partial \rho}{\partial t} = 0$$
\end{theorem}

\textbf{Interprétation :} Si la densité de charge $\rho$ diminue en un point, c'est qu'un courant $\vec{j}$ s'en éloigne (divergence positive).

\section{Les Équations de Maxwell}

\subsection{Équations de structure}
Elles décrivent les propriétés intrinsèques des champs $\vec{E}$ et $\vec{B}$.

\begin{definition}[Maxwell-Gauss et Maxwell-Thomson]
    \begin{itemize}
        \item \textbf{Maxwell-Gauss :} $\Div \vec{E} = \frac{\rho}{\epsilon_0}$ (Lien source-champ)
        \item \textbf{Maxwell-Thomson :} $\Div \vec{B} = 0$ (Flux conservatif, pas de monopôle magnétique)
    \end{itemize}
\end{definition}

\subsection{Équations de couplage (régime variable)}
Ces équations traduisent l'interaction entre les champs et leur dépendance temporelle.

\begin{theorem}[Maxwell-Faraday et Maxwell-Ampère]
    \begin{itemize}
        \item \textbf{Maxwell-Faraday :} $\rot \vec{E} = -\frac{\partial \vec{B}}{\partial t}$ 
        \item \textbf{Maxwell-Ampère :} $\rot \vec{B} = \mu_0\left(\vec{j} + \epsilon_0\frac{\partial \vec{E}}{\partial t}\right)$
    \end{itemize}
\end{theorem}

\tipsInfo{Courant de déplacement}{Le terme $\vec{j_D} = \epsilon_0\frac{\partial \vec{E}}{\partial t}$ est appelé \textbf{courant de déplacement}. Il a été introduit par Maxwell pour assurer la cohérence mathématique avec l'équation de conservation de la charge.}

\section{Potentiels Électromagnétiques}

On définit les potentiels $(V, \vec{A})$ qui permettent de reconstruire les champs physiques.

\begin{definition}[Potentiels vecteur et scalaire]
    Les champs $\vec{E}$ et $\vec{B}$ s'expriment en fonction du potentiel vecteur $\vec{A}$ et du potentiel scalaire $V$ par :
    $$\vec{B} = \rot \vec{A} \quad \text{et} \quad \vec{E} = -\grad V - \frac{\partial \vec{A}}{\partial t}$$
\end{definition}

\subsection{Jauge de Lorenz}
Pour simplifier les équations de propagation, on exploite la liberté de jauge.

\begin{theorem}[Condition de Lorenz]
    On impose la relation suivante entre $\vec{A}$ et $V$ :
    $$\Div \vec{A} + \frac{1}{c^2}\frac{\partial V}{\partial t} = 0 \quad \text{avec} \quad c = \frac{1}{\sqrt{\mu_0\epsilon_0}}$$
\end{theorem}

\subsection{Équations de propagation}


Sous la jauge de Lorenz, les potentiels sont découplés et satisfont des équations d'onde de type D'Alembertien.

\begin{theorem}[Équations de propagation des potentiels]
    $$\square V = \Delta V - \frac{1}{c^2}\frac{\partial^2 V}{\partial t^2} = -\frac{\rho}{\epsilon_0}$$
    $$\square \vec{A} = \Delta \vec{A} - \frac{1}{c^2}\frac{\partial^2 \vec{A}}{\partial t^2} = -\mu_0 \vec{j}$$
\end{theorem}

\subsection{Invariance de jauge}
Il existe une infinité de couples $(V, \vec{A})$ pour décrire les mêmes champs physiques.

\begin{theorem}[Transformation de jauge]
    Les champs $\vec{E}$ et $\vec{B}$ sont invariants par la transformation :
    $$\vec{A}' = \vec{A} + \grad f \quad ; \quad V' = V - \frac{\partial f}{\partial t}$$
    où $f(M,t)$ est un champ scalaire quelconque.
\end{theorem}

\end{document}